\chapter{Hybrid compressed and addressable memory}
\label{evod40:chapter}

A source event says ``alpha beta omitted-x.'' Another says ``alpha beta
omitted-y.'' A compressor retaining the first two distinct words gives both
the same text. Can the agent later tell which event occurred? A larger context
window does not answer this question. The system must preserve the source
record, its identity and the link from summary to source---or explicitly accept
that the distinction is gone.

Chapter 20 derived compressed numerical memory and an order-loss
counterexample. Chapter 35 built token-addressable Transformer execution
state. Chapter 39 coordinated those engines with retrieval and tools.
Here we add a separate memory contract: durable events, lossy derived records,
addressable reads, bounded working context and explicit unavailability.
The reference is a serialized in-memory specification of those rules.
Calling its objects durable names their intended lifecycle; it does not make
this Python dictionary survive a process crash.

\section{Four memory classes answer different questions}
\label{evod40:classes}
\index{memory!native versus durable}
Engine-native transient state supports numerical continuation. For the
restricted native recurrence it includes hidden vector and token position;
for Hermon it includes every layer's ordered KV and model/position metadata.
A checkpoint may persist those bytes, but that does not make them searchable
semantic evidence.

An immutable \Term{durable event} records an occurrence under a tenant and
schema. A \Term{summary} is a lossy derived record with explicit source edges.
An \Term{addressable record} has a stable identity that a lexical, vector,
graph or temporal index can return. Addressability is an access property;
events and summaries can both be addressable. These are not four disjoint
physical storage devices.

A native checkpoint can be useful for exact continuation but poor for
semantic search. A document can answer a retrieval query yet be insufficient
to reconstruct the exact numerical state. A system may keep both, linked
through provenance and artifact identities. It must not serialize a retrieved
paragraph into a state tensor and claim exact continuation without a tested
decoder relation.

\section{Define identity, time and availability before ranking}
\label{evod40:records}
\index{record identity}\index{expiry}\index{provenance!ancestor availability}
Our record is
\[
 R=(id,tenant,kind,payload,sources,t_c,t_e,artifact).
\]
$t_c$ and optional $t_e$ count \emph{logical integer ticks}, not wall-clock
seconds. A record expires at $t\ge t_e$, including equality. An event's
artifact names its schema; a summary's artifact names the compressor and word
limit. Production records would also pin data/prompt/policy versions, creation
authority, quality evidence and content encodings.

The ID hashes the record fields and a fresh store serial using SHA-256.
Including serial prevents a forgotten event from being resurrected by
reinserting identical text under the same ID. An explicit later event with
equal payload is a new occurrence; it does not restore old provenance edges.
These hashes do not authenticate producers, protect a hostile object store
or replace a transactional ID service.

Define $L_t(r)$ as the conjunction of existence, tenant match, absence of a
tombstone and non-expiry. For an event, availability is $A_t(r)=L_t(r)$.
For a derived record,
\begin{equation}
 A_t(r)=L_t(r)\ \wedge\!\bigwedge_{s\in sources(r)} A_t(s).
 \label{evod40:availability}
\end{equation}
This is a strict evidence policy: if any ancestor is unavailable, the summary
is unavailable. Another system might return a stale summary with explicit
warnings, but silently presenting it as current evidence is a different
contract. Read-time validation dominates index contents.

\Listing{Memory._available}
Summary creation accepts only existing available distinct source IDs.
Edges therefore point to older registered objects, forming a DAG. The store
bounds provenance depth to 32 and total records to 10,000; it also detects
cycles if trusted internal storage is corrupted. The recursive checker memoizes
shared ancestors within one availability call, bounding it by the reachable
vertices and edges. Retrieval repeats that check for candidate records.
A production DAG still needs incremental invalidation and deletion-transaction
design; no scalable lookup claim follows from this small reference.

\Plate{01-provenance}{Deletion propagates through evidence dependencies. The surviving source remains readable, but both derived records fail the strict ancestor-availability rule.}{evod40:dag}

\par\medskip\noindent\begin{minipage}{\linewidth}
\begin{proposition}[Deletion dominance under fixed records]
Assume a finite acyclic source graph, fixed record bytes, nondecreasing time,
and tombstones that are never removed. Once a record is unavailable because
of expiry or a tombstoned ancestor, no later read of that record becomes
available merely by updating or retaining an index entry.
\end{proposition}
\begin{proof}
For a source event, neither time advancing nor adding tombstones can turn
$L_t$ from false to true. Induct in source-before-derived order. Every
derived availability is a conjunction of its own local predicate and already
monotone ancestor predicates. That conjunction cannot turn true after one
such cause makes it false. Ranking does not appear in
Equation~\ref{evod40:availability}, so stale index entries cannot alter it.
\end{proof}
\end{minipage}\par\medskip
The fixed-record premise matters. A corrected event is a new record and
requires new summaries. Editing bytes in place would invalidate both identity
and this reasoning.

\section{Compression preserves provenance, not all information}
\label{evod40:compression}
\index{compression!non-injectivity}\index{memory!payload accounting}
The supplied package's compressor keeps the first distinct words encountered
over ordered sources. We retain that transparent rule because it makes loss
inspectable, rather than pretending an untrained summarizer has learned
semantic retention. Source IDs remain ordered and the compressor version
includes the word limit.
\Listing{Memory.summarize}

\Plate{02-alias}{A payload collision is an exact information-loss witness. Provenance can distinguish the records even when their summary texts are identical.}{evod40:alias}

For $C(X)=Z$, if $C(X_1)=C(X_2)$ for distinct source histories, no function
of $Z$ alone can always recover the history: the same input $Z$ would require
two outputs. This elementary non-injectivity argument is independent of
floating-point precision. Source expansion is not a general inverse of $C$;
it is an addressable read of separately retained evidence.

Measure compression and usefulness separately:
\[
 \rho_{payload}=\frac{\operatorname{bytes}_{UTF8}(Z)}
 {\sum_i\operatorname{bytes}_{UTF8}(X_i)},\qquad
 U=\operatorname{quality}(task\mid Z).
\]
The generated fixture stores the exact numerator and denominator in
\texttt{results.json}. This ratio excludes IDs, source edges, object overhead,
index entries and retained original records. Keeping the originals \emph{and}
adding a summary increases total stored bytes even if its payload ratio is
small. A summary useful for approximate topic recall can fail exact quotation,
chronology or a rare identifier query. No utility measurement is implemented
by counting words.

\section{Rank records only after scope and provenance checks}
\label{evod40:retrieval}
\index{retrieval!lexical ranking}\index{index!stale entries}
The addressable query returns ordered record IDs with payloads after tenant,
retention and ancestor checks. The reference lexical score is the number of
distinct lowercase ASCII alphanumeric query terms also present in a record.
Scores are nonnegative integers, not probabilities or semantic similarity
measurements. Zero-overlap records are omitted. Ties use creation tick and
full record ID, making the result deterministic within a fixed store image.

An optional candidate-ID list models a stale index. Before returning a hit,
the object store rechecks current availability. A tombstoned ID still present
in that list cannot bypass deletion; duplicates are removed. Cross-tenant
candidate IDs produce no hits. Direct foreign-object reads raise a permission
error in this teaching contract. Revealing existence through error differences
can itself require a different policy in a hostile multi-tenant service.

Retrieval should be conditional on an evidence need, not blindly invoked
before every token. MemGPT motivates explicit context tiers and movement
between them~\cite{evod40:memgpt}. Self-RAG explicitly makes retrieval
adaptive~\cite{evod40:selfrag}. Those narrow precedents do not establish that
this lexical ranker reproduces either system, nor that a retrieved result is
correct or safe. Their reported benchmark comparisons are not our evidence.

\section{Bound source expansion and working-context promotion}
\label{evod40:rehydrate}
\index{Rehydration!budgets}\index{UTF-8 payload bytes}
For summary $s$, rehydration traverses source links breadth-first, visiting
each ID at most once in its returned expansion. It accepts both a record-count
ceiling $B_n$ and a UTF-8 payload-byte ceiling $B_b$. The first whole record
that would exceed the remaining byte budget stops traversal. This deterministic
prefix policy favors provenance order rather than solving a knapsack problem.
\Listing{Memory.rehydrate}

If returned records are $r_1,\ldots,r_m$, the maintained invariant is
\[
 m\le B_n,\qquad used=\sum_{j=1}^m |payload(r_j)|_{UTF8}\le B_b.
\]
Initially both sums are zero. A record is appended only if its full byte size
fits; then both count and sum advance together. The loop guard enforces the
count. This proves the payload bound by induction, not a bound on Python
object memory or tokenizer context length. Repeated paths to a shared source
do not duplicate that source in the result.

Promotion selects whole readable records into a \Term{working set} under a
declared byte ceiling. Unlike rehydration's prefix, the direct promotion
operation is all-or-error: if the requested set is too large, it returns no
partial context. Both interfaces are explicit so callers can choose rather
than infer an accidental truncation rule.
\Listing{Memory.promote}
\Plate{03-budget}{Payload budgets limit context movement. Authorization comes from the separate policy configuration, never from the selected record text.}{evod40:budget}

UTF-8 bytes are not tokens. Two Chinese characters consume six UTF-8 bytes
before spaces or other text; their tokenizer length depends on the exact
tokenizer artifact. A real planner must reserve engine-context tokens after
the declared conversion, as well as storage and transport bytes.

\section{Retention, forgetting and the evidence firewall}
\label{evod40:firewall}
\index{forgetting!issued contexts}\index{authority!memory firewall}
Expiry expresses a declared availability horizon. Forgetting adds a tombstone
even for an already expired record and is idempotent. It does not delete the
dictionary payload. Retention policy may additionally require physical erasure,
backup expiry, legal holds, index compaction or derived-data destruction;
those actions are not implemented here.

An already returned immutable context still contains its payload after a
source is forgotten. The store can reject future promotion, but it cannot
recall bytes already issued to a client or a running model. Strong revocation
therefore needs request-time fences, context leases or cancellation and an
explicit policy for in-flight work. Claiming instantaneous distributed erasure
from a local tombstone test would be false.

The \Term{memory safety firewall} is a semantic authority boundary. Every
\texttt{Context} has an empty authority tuple; memory operators cannot mint
Chapter 39 policy grants. Retrieved text can contain stale facts, secrets or
prompt injection. Provenance and quoting let a planner reason about it as
evidence. Authorization remains an independent configured decision before
an action. This small typed separation is not a production prompt-injection
defense or proof of hostile-process isolation.

\section{Consistency and failure controls}
\label{evod40:evidence}
\index{consistency!memory store}
\begin{center}\input{\Generated/visibility.tex}\end{center}
At tick zero, both source events and the derived summary can contribute
ranked hits. At tick two, one source expires; its summary becomes unavailable
under the strict rule while the unrelated live event remains retrievable.
The exact compression collision has equal summary text and unequal source IDs.
These are deterministic logical examples, not retrieval-quality benchmarks.

Tests cover exact TTL equality, transitive invalidation, stale index entries,
cross-tenant reads, duplicate sources, missing provenance, shared ancestors,
UTF-8 budgets, corrupted cycles and equal-payload reinsertion. A negative
control deliberately retains an issued context after deletion to demonstrate
the revocation limit. Testing only direct event lookup would miss the stale
summary path.

The reference assumes serialized trusted callers and strongly consistent
dictionary access. A real service must define atomicity among object writes,
source edges, indexes and tombstones; conflict resolution; replication lag;
deletion epochs; and failure recovery. A safe eventual index can lag while
the authoritative read gate stays current. A stale authoritative tombstone
replica needs a separate consistency protocol. Chapter 42's recovery model
will provide vocabulary, not automatically implement that distributed store.

\section{Twelve worked exercises}
\label{evod40:exercises}
\begin{enumerate}
\item Is KV automatically durable memory? \textbf{Answer:} No. It supports
native continuation. Persistent searchable evidence needs identity, retention
and provenance beyond a transient cache.
\item Show a compression collision. \textbf{Answer:} The two strings ending
in omitted-x and omitted-y both map to alpha beta under the two-word rule.
\item Why is rehydration not inversion? \textbf{Answer:} It reads retained
source IDs. If those objects are unavailable, the summary cannot recreate them.
\item An event expires at tick 2. Is it readable at tick 2?
\textbf{Answer:} No; the availability predicate requires $t<t_e$.
\item A summary depends on a forgotten source through another summary.
What happens? \textbf{Answer:} Induction through the provenance DAG makes
both derived records unavailable, even if indexes list them.
\item Sources have payload sizes 10, 10 and 12 bytes. What does prefix
rehydration return under 20 bytes and three records? \textbf{Answer:} The
first two, totaling 20 bytes; it does not cut the third record.
\item Does a 20-byte bound imply 20 tokenizer tokens? \textbf{Answer:} No.
The byte encoding and tokenizer artifact define different units.
\item What happens when identical forgotten text is inserted again?
\textbf{Answer:} It gets a fresh occurrence ID. Old summaries still point
to the old tombstoned source and remain unavailable.
\item Can an already returned context disappear after forgetting?
\textbf{Answer:} Not from this caller's memory. Future reads can be denied;
revoking in-flight use needs a separate lifecycle protocol.
\item Why distinguish payload ratio from total storage?
\textbf{Answer:} The latter includes original records, links, identities and
index overhead. Summaries added to retained sources can increase storage.
\item How would you evaluate a learned compressor? \textbf{Rubric:} Pin
source sets, compressor/prompt and budgets; compare raw records, truncation
and summary-plus-retrieval on held-out exact and approximate tasks; include
rare identifiers, order changes, deleted sources and adversarial instructions;
report utility and full memory cost separately.
\item What must a distributed memory contract specify?
\textbf{Answer:} Which store is authoritative, visibility/transaction boundaries,
replication and conflict rules, tombstone propagation, recovery and in-flight
revocation. Naming a memory tier supplies none of those guarantees.
\end{enumerate}

\section{Vocabulary and the compiler boundary}
\Term{Provenance DAG}: source relations among immutable evidence records.
\Term{Tombstone}: an explicit denial of future access, not proof of physical
erasure. \Term{Rehydration}: bounded source traversal rather than loss reversal.
Chapter 41 now treats engine, memory, retrieval and tool operations as typed
instructions whose effects must survive legal compilation.
